
% Latex source to make a diagram of the UM subroutine call tree, showing the
% locations of all cloud scheme calls.  This diagram is included in both
% UMDP 029 (large-scale cloud scheme) and UMDP 030 (PC2).

% NOTE: any preamble text required for this should be put in the file
% um_call_tree_preamble.tex, which is also inlcuded in both the UMDPs.

Subroutines only called for the \textcolor{blue}{Smith} scheme are highlighted
in \textcolor{blue}{blue}, those only called for \textcolor{mygreen}{PC2} are
in \textcolor{mygreen}{green}, and those only called for
the \textcolor{purple}{bimodal} scheme are in \textcolor{purple}{purple}.

\subsubsection{Main Tree from atm\_step\_4a}

\begin{itemize}

\item {\bf atm\_step\_4a} \\*
(performs one timestep of the Unified Model...)
  \begin{itemize}

  \begin{tcolorbox}[enhanced jigsaw, breakable]
  \item {\bf atm\_step\_alloc\_4a} \\*
  (does miscellaneous initialisations in atm\_step)

    \begin{itemize}

    \item \textcolor{mygreen}{\bf pc2\_rhtl} \\*
    (calculate start-of-timestep Relative Humidity, used by PC2 initiation)

    \end{itemize}
  \end{tcolorbox}

  \begin{tcolorbox}[enhanced jigsaw, breakable]
  \item {\bf atmos\_physics1} \\*
  (calls explicit ``slow'' physics routines...)
    \begin{itemize}

    \begin{tcolorbox}
    \item {\bf microphys\_ctl} \\*
    (interface to microphysics scheme)
      \begin{itemize}

      \item \textcolor{mygreen}{\bf pc2\_turbulence\_ctl} \\*
      (Perform optional erosion of liquid-cloud;
       done here if NOT doing erosion after convection,
       e.g. if no convection scheme is used).
        \begin{itemize}

        \item \textcolor{mygreen}{\bf pc2\_hom\_conv} \\*
        (called here just to do erosion)

        \end{itemize}

      \item \textcolor{blue}{\bf ls\_cld} \\*
      (Smith scheme without area cloud fraction calculation,
       to set initial cloud fields passed into microphysics)

      \item {\bf ls\_ppn} \\*
      (microphysics scheme)

      \item {\bf mphys\_turb\_gen\_mixed\_phase} \\*
      (turbulent production of liquid cloud)

      \item \textcolor{mygreen}{\bf pc2\_turbulence\_ctl} \\*
      (optionally use the PC2 pdf-width-change code to calculate
       the cloud-fraction change from the above turbulent production
       of liquid cloud)

        \begin{itemize}

        \item \textcolor{mygreen}{\bf pc2\_hom\_conv} \\*
        (called here just to calculate the cloud fraction increment
         consistent with the turbulent qcl increment)

        \end{itemize}

      \end{itemize}
    \end{tcolorbox}

    \begin{tcolorbox}
    \item {\bf rad\_ctl} \\*
    (interface to radiation scheme)
      \begin{itemize}

      \item {\bf sw\_rad} \\*
      (short-wave radiation scheme)

      \item \textcolor{mygreen}{\bf pc2\_homog\_plus\_turb} \\*
      (PC2 homogeneous forcing of liquid-cloud by SW radiation heating)

      \item {\bf lw\_rad} \\*
      (long-wave radiation scheme)

      \item \textcolor{mygreen}{\bf pc2\_homog\_plus\_turb} \\*
      (PC2 homogeneous forcing of liquid-cloud by LW radiation tendency)

      \end{itemize}
    \end{tcolorbox}

    \begin{tcolorbox}
    \item {\bf atmos\_physics1\_alloc\_pc2}
    (wrapper for PC2 self-consistency checks at end of atmos\_physics1)

      \begin{itemize}

      \item Add increments from microphysics + radiation onto
      start-of-timestep fields to form updated fields.

      \item \textcolor{mygreen}{\bf pc2\_checks} \\*
      (self-consistency checks on cloud fractions and water contents)

      \item Convert corrected updated fields back to increments.

      \end{itemize}
    \end{tcolorbox}

    \end{itemize}
  \end{tcolorbox}

  Begin loop over solver outer cycles

    \begin{itemize}

    \begin{tcolorbox}[enhanced jigsaw, breakable]
    \item {\bf atm\_step\_phys\_reset} \\*
    (for PC2, on subsequent solver outer cycles,
     reset cloud-fractions to saved values after atmos\_physics1)
    \end{tcolorbox}

    \begin{tcolorbox}[enhanced jigsaw, breakable]
    \item {\bf eg\_sl\_moisture} \\*
    (large-scale advection of cloud water contents and fractions)
    \end{tcolorbox}

    \begin{tcolorbox}[enhanced jigsaw, breakable]
    \item \textcolor{mygreen}{\bf pc2\_pressure\_forcing\_only} \\*
    (Optionally calculate homogeneous forcing of liquid cloud by the
     pressure change along the trajectory from departure point to
     arrival point).
       \begin{itemize}

       \item \textcolor{mygreen}{\bf pc2\_homog\_plus\_turb} \\*
       (generic homogeneous forcing routine used here).

       \end{itemize}
    \end{tcolorbox}

    \begin{tcolorbox}[enhanced jigsaw, breakable]
    \item {\bf atmos\_physics2} \\*
    (calls ``fast'' physics routines...)

      \begin{itemize}

      \begin{tcolorbox}
      \item {\bf ni\_bl\_ctl} \\*
      (interface to explicit boundary-layer and surface scheme calls,
       including calculation of TKE and TKE-based $RH_{crit}$)
      \end{tcolorbox}

      \begin{tcolorbox}
      \item \textcolor{purple}{\bf bm\_calc\_tau} \\*
      (calculates turbulence properties used in the bimodal cloud scheme,
       based on the boundary-layer scheme TKE and mixing-length)
      \end{tcolorbox}

      \begin{tcolorbox}
      \item {\bf cloud\_call\_b4\_conv} \\*
      (routine for optional cloud-scheme calls before convection)
        \begin{itemize}

        \item \textcolor{blue}{\bf ls\_arcld} \\*
        (Smith scheme with area cloud fraction;
         see \ref{subsubsec:smith_acf} for a drill-down inside this routine)

        \item \textcolor{purple}{\bf bm\_ctl} \\*
        (bimodal scheme)
    
        \item \textcolor{purple}{Set area cloud fraction equal to bulk
                                 cloud fraction}

        \item \textcolor{mygreen}{\bf pc2\_initiation\_ctl} \\*
        (interface to PC2 initiation and consistency-checks;
         see \ref{subsubsec:pc2_initiation} for a drill-down inside this
         routine)

        \end{itemize}
      \end{tcolorbox}

      \begin{tcolorbox}
      \item {\bf ni\_conv\_ctl} or {\bf other\_conv\_ctl} \\*
      (interface routines to various convection schemes...)
        \begin{itemize}

        \item {\bf glue\_conv\_5a/6a} \\*
        (calls deep, shallow and mid-level convection schemes)
          \begin{itemize}

          \item{\bf deep/shallow/mid\_conv} \\*
          (convection scheme main routines)
            \begin{itemize}

            \item {\bf convec2}
            (completes lifting of the convective parcel by one model-level)
              \begin{itemize}

              \item {\bf parcel}
              (calculates new parcel properties at next level)

              \item {\bf environ}
              (calculates grid-mean increments to primary fields;
               includes PC2 partitioning of detrained condensate mass
               between liquid and ice phases)

              \item \textcolor{mygreen}{\bf pc2\_environ}
              (calculates increments to PC2 cloud fractions due to
               convective detrainment and subsidence)

              \end{itemize}

            \end{itemize}

          \end{itemize}

        \item \textcolor{mygreen}{\bf pc2\_from\_conv\_ctl} \\*
        (PC2 calculations after convection)
          \begin{itemize}

          \item \textcolor{mygreen}{\bf pc2\_hom\_conv} \\*
          (homogeneous forcing by convection, and erosion of liquid-cloud)

          \end{itemize}

        \end{itemize}
      \end{tcolorbox}

      \begin{tcolorbox}
      \item {\bf ni\_imp\_ctl} \\*
      (interface to boundary-layer implicit solver)
        \begin{itemize}

        \item {\bf imp\_solver} \\*
        (implicitly solves vertical diffusion to find $T_l$ and $q_T$
         updated by turbulent fluxes).

        \item \textcolor{mygreen}{\bf pc2\_bl\_inhom\_ice} \\*
        (inhomogeneous forcing of ice-cloud)

        \item \textcolor{mygreen}{\bf pc2\_delta\_hom\_turb} \\*
        (homogeneous forcing of liquid cloud by the turbulent fluxes)

        \item \textcolor{mygreen}{\bf pc2\_bl\_forced\_cu} \\*
        (adds diagnosed ``forced cumulus'' cloud fraction and water content
         onto the PC2 prognostics)

        \item Calculate area cloud fraction:

        \textcolor{mygreen}{\bf ls\_acf\_brooks} \\*
        (for the Brooks epirical method)

        \textcolor{mygreen}{\bf pc2\_hom\_arcld} \\*
        (for the Cusack vertical interpolation method)
          \begin{itemize}

          \item \textcolor{mygreen}{\bf pc2\_homog\_plus\_turb} \\*
          (generic homogeneous forcing routine used to interpolate)

          \end{itemize}

        \item \textcolor{blue}{\bf ls\_arcld} \\*
        (interface to diagnostic Smith scheme and area cloud fraction;
         see \ref{subsubsec:smith_acf} for a drill-down inside this routine)

        \item \textcolor{purple}{\bf bm\_ctl} \\*
        (bimodal cloud scheme)

        \item \textcolor{purple}{Set area cloud fraction equal to bulk
                                 cloud fraction}

        \item {\bf diagnostics\_bl} \\*
        (outputs boundary-layer diagnostics to STASH)
          \begin{itemize}

          \item {\bf ls\_cld} \\*
          (Smith scheme used here to calculate various diagnostics of
           near-surface temperature and humidity, by extrapolating pressure,
           $T_l$ and $q_t$ down to the desired height and then
           re-diagnosing $q_{cl}$.

          \end{itemize}

        \end{itemize}
      \end{tcolorbox}

      \end{itemize}
    \end{tcolorbox}

    \begin{tcolorbox}[enhanced jigsaw, breakable]
    \item {\bf atm\_step\_ac\_assim} \\*
    (interface to Data Assimilation analysis increments...)

      \begin{itemize}

      \item {\bf ac\_ctl}
      (control routine for Data Assimilation analysis increments...)
        \begin{itemize}

        \item{\bf ac}
        (main analysis increment routine)

        \item \textcolor{mygreen}{\bf pc2\_assim} \\*
        (PC2 reponse to the analysis increments;
         see \ref{subsubsec:pc2_assim} for a drill-down inside this routine)

        \item \textcolor{mygreen}{\bf ls\_acf\_brooks}
        (calculate area cloud fraction using Brooks empirical method if active)

        \item \textcolor{blue}{\bf ls\_arcld}
        (call diagnostic Smith scheme with area cloud fraction again to
         account for the analysis increments;
         see \ref{subsubsec:smith_acf} for a drill-down inside this routine)

        \end{itemize}

      \end{itemize}
    \end{tcolorbox}

    \begin{tcolorbox}[enhanced jigsaw, breakable]
    \item {\bf eg\_sl\_helmholtz} \\*
    (dynamics pressure solver; updates pressure, and the winds used
     to perform advection on the next solver outer cycle)
    \end{tcolorbox}

    \end{itemize}

  End loop over solver outer cycles

  \begin{tcolorbox}[enhanced jigsaw, breakable]
  \item \textcolor{mygreen}{\bf pc2\_pressure\_forcing} \\*
  (interface to miscellaneous PC2 calculations at end-of-timestep)
    \begin{itemize}

    \item \textcolor{mygreen}{\bf pc2\_homog\_plus\_turb} \\*
    (homogeneous forcing of liquid-cloud by the dynamics pressure change;
     optionally either uses total pressure change including the
     Lagrangian component following the winds, or only the Eulerian
     component from the dynamics solver)

    \item \textcolor{mygreen}{\bf pc2\_initiation\_ctl} \\*
    (interface to PC2 initiation and consistency-checks;
     see \ref{subsubsec:pc2_initiation} for a drill-down inside this routine)

    \end{itemize}
  \end{tcolorbox}

  \begin{tcolorbox}[enhanced jigsaw, breakable]
  \item {\bf qt\_bal\_cld} \\*
  (calculates end-of-timestep cloud state consistent with final pressure...)
    \begin{itemize}

    \item \textcolor{blue}{\bf ls\_arcld} \\*
    (interface to diagnostic Smith scheme and area cloud fraction;
     see \ref{subsubsec:smith_acf} for a drill-down inside this routine)

    \item \textcolor{purple}{\bf bm\_ctl} \\*
    (bimodal cloud scheme)

    \item \textcolor{purple}{Set area cloud fraction equal to bulk
                             cloud fraction}

    \end{itemize}
  \end{tcolorbox}

  \begin{tcolorbox}[enhanced jigsaw, breakable]
  \item {\bf iau} \\*
  (incremental analysis update; part of data assimilation)
    \begin{itemize}

    \item \textcolor{mygreen}{\bf pc2\_assim} \\*
    (PC2 reponse to the analysis increments;
     see \ref{subsubsec:pc2_assim} for a drill-down inside this routine)

    \item \textcolor{mygreen}{\bf initial\_pc2\_check} \\*
    (wrapper for optional self-consistency checks on prognostic cloud variables
     if not doing PC2 response to analysis increments)

      \begin{itemize}

      \item \textcolor{mygreen}{\bf pc2\_checks} \\*
      (self-consistency checks on cloud fractions and water contents)

      \end{itemize}

    \end{itemize}
  \end{tcolorbox}

  \end{itemize}

\end{itemize}


Drill-downs within some routines in the call tree are listed separately below,
to avoid duplication
(since these routines are called in multiple different places in the tree)...

\subsubsection{Smith scheme with area cloud fraction}
\label{subsubsec:smith_acf}

\begin{itemize}

\begin{tcolorbox}[enhanced jigsaw, breakable]
\item \textcolor{blue}{\bf ls\_arcld} \\*
(interface to diagnostic Smith scheme and area cloud fraction)
  \begin{itemize}

  \item If no area cloud fraction scheme:

  \textcolor{blue}{\bf ls\_cld} \\*
  (just directly call Smith scheme)

  Set area cloud fraction equal to bulk cloud fraction.

  \item If using Cusack vertical interpolation method:

  Interpolate fields onto finer vertical grid

  \textcolor{blue}{\bf ls\_cld} \\*
  (call Smith scheme using higher vertical resolution fields)

  Coarse-grain cloud fields back to model grid, but set area cloud
  fraction to max of bulk cloud fraction over corresponding fine-grid levels.

  \item If using Brooks empirical area cloud fraction method:

  \textcolor{blue}{\bf ls\_cld} \\*
  (just directly call Smith scheme)

  \textcolor{blue}{\bf ls\_acf\_brooks} \\*
  (estimate area cloud fraction)

  \end{itemize}
\end{tcolorbox}

\end{itemize}


\subsubsection{PC2 initiation}
\label{subsubsec:pc2_initiation}

\begin{itemize}

\begin{tcolorbox}[enhanced jigsaw, breakable]
\item \textcolor{mygreen}{\bf pc2\_initiation\_ctl} \\*
(interface to PC2 initiation and consistency-checks)
  \begin{itemize}

  \item \textcolor{mygreen}{\bf pc2\_checks} \\*
  (self-consistency checks on cloud fractions and water contents)

  \item PC2 initiation of liquid-cloud:

  \textcolor{mygreen}{\bf pc2\_bm\_initiate} \\*
  (using the bimodal cloud scheme)

  \textcolor{mygreen}{\bf pc2\_arcld} \\*
  (using the Smith scheme with the Cusack vertical interpolation method)

    \begin{itemize}

    \item \textcolor{mygreen}{\bf pc2\_initiate} \\*
    (initiation using the Smith scheme,
     called here on a finer vertical grid as per the Cusack method)

    \end{itemize}

  \textcolor{mygreen}{\bf pc2\_initiate} \\*
  (using the Smith scheme with no area cloud representation)

  \item \textcolor{mygreen}{\bf pc2\_checks2} \\*
  (further self-consistency checks on cloud-fractions)

  \item \textcolor{mygreen}{\bf pc2\_checks} \\*
  (repeat the first lot of self-consistency checks again,
   just in case we broke something in the mean-time!)

  \item \textcolor{mygreen}{\bf pc2\_hom\_arcld} \\*
  (finds area cloud fraction using a version of the Cusack method,
   where the cloud fraction on the finer vertical grid is estimated by
   applying homogeneous forcing relative to the original grid fields)

    \begin{itemize}

    \item \textcolor{mygreen}{\bf pc2\_homog\_plus\_turb} \\*
    (generic homogeneous forcing routine used to interpolate)

    \end{itemize}

  \end{itemize}
\end{tcolorbox}

\end{itemize}


\subsubsection{PC2 Data Assimilation}
\label{subsubsec:pc2_assim}

\begin{itemize}

\begin{tcolorbox}[enhanced jigsaw, breakable]
\item \textcolor{mygreen}{\bf pc2\_assim} \\*
(PC2 reponse to the analysis increments)
  \begin{itemize}

  \item \textcolor{mygreen}{\bf pc2\_homog\_plus\_turb} \\*
  (generic PC2 homogeneous forcing routine used here for liquid-cloud)

  \item Estimate change in ice-cloud fraction from the assimilation
  increment to ice-cloud mass.

  \item \textcolor{mygreen}{\bf pc2\_total\_cf} \\*
  (update bulk cloud fraction due to change in ice cloud fraction)

  \item \textcolor{mygreen}{\bf pc2\_checks} \\*
  (self-consistency checks on prognostic cloud fractions and
   water contents)

  \end{itemize}
\end{tcolorbox}

\end{itemize}
